%% ma: L12-B08-0186
%% lop: 12
%% chuong: 2
%% bai: 8
%% dang: 12.8.4
%% loai: TN4
%% muc: VD
%% dap_an: "A"
%% nguon: "CĐ2.6.pdf giáo viên gửi 27/10/2026 (M07, Lê Minh Khôi – Nguyễn Anh Khôi, Chương II Chủ đề 6 Ứng dụng vào các bài toán thực tế), Câu 14"
%% kien_thuc: "Bài 8 (tọa độ vectơ, tích vectơ với một số, khoảng cách), chuyển động thẳng đều"
%% trang_thai: chinh-thuc
%% ghi_chu: "Đề gốc dùng chung giả thiết và hình cho Câu 13–14; đã chép vào từng câu"
\begin{ex}
Ba chiếc flycam I, II, III được điều khiển cùng bay lên tại một địa điểm. Sau một thời gian bay, chiếc flycam thứ nhất (I) cách mặt đất $3$ m, cách điểm xuất phát $3$ m về phía Nam và $2$ m về phía Đông. Chiếc flycam thứ hai (II) cách mặt đất $6$ m, cách điểm xuất phát $9$ m về phía Nam và $9$ m về phía Đông. Chiếc flycam thứ ba (III) cách mặt đất $6$ m, cách điểm xuất phát $6$ m về phía Bắc và $4$ m về phía Tây. Quy ước hệ trục tọa độ $Oxyz$ với gốc tọa độ $O$ trùng với điểm xuất phát của ba chiếc flycam, mặt đất là tập hợp các điểm có cao độ bằng $0$, hai tia $Ox$, $Oy$ lần lượt hướng về phía Nam và phía Đông, đơn vị dài trên mỗi trục là mét (hình minh họa).
\begin{center}
\begin{tikzpicture}[scale=.3,x={(-.4cm,-.5cm)},y={(1cm,0cm)},z={(0cm,1cm)}]
  \fill[green!10] (12,-8,0)--(12,12,0)--(-8,12,0)--(-8,-8,0)--cycle;
  \draw[truc] (0,0,0)--(13,0,0) node[below left]{$x$ (Nam)};
  \draw[truc] (0,-8,0)--(0,13,0) node[right]{$y$ (Đông)};
  \draw[truc] (0,0,0)--(0,0,9) node[right]{$z$};
  \path (3,2,3) coordinate (I) (9,9,6) coordinate (II) (-6,-4,6) coordinate (III);
  \draw[phu] (3,0,0)--(3,2,0)--(0,2,0) (3,2,0)--(I);
  \draw[phu] (9,0,0)--(9,9,0)--(0,9,0) (9,9,0)--(II);
  \draw[phu] (-6,0,0)--(-6,-4,0)--(0,-4,0) (-6,-4,0)--(III);
  \foreach \P in {I,II,III} \fill (\P) circle (6pt);
  \path (I) node[right]{I} (II) node[right]{II} (III) node[left]{III} (0,0,0) node[below right]{$O$};
\end{tikzpicture}
\end{center}
Giả sử ba chiếc flycam I, II, III bay thẳng đều với cùng tốc độ $v$ (m/s) từ điểm xuất phát đến các vị trí đã nêu. Biết chiếc flycam đến đích đầu tiên thì ngay lúc đó hai chiếc còn lại đổi hướng bay thẳng đến đích của chiếc đầu tiên. Tổng quãng đường bay của hai chiếc flycam kể từ lúc đổi hướng gần nhất với giá trị nào sau đây?
\choice{\True $8{,}63$ m}{$9{,}15$ m}{$10{,}20$ m}{$7{,}85$ m}
\loigiai{$OI=\sqrt{22}$, $OII=\sqrt{198}=3\sqrt{22}$, $OIII=\sqrt{88}=2\sqrt{22}$: flycam I đến đích trước.

Khi đó II đi được $\dfrac13$ quãng đường, ở $(3;3;2)$; III đi được $\dfrac12$, ở $(-3;-2;3)$.

Tổng quãng đường sau khi đổi hướng: $\sqrt{0+1+1}+\sqrt{36+16+0}=\sqrt2+\sqrt{52}\approx8{,}63$ m.}
\end{ex}
